\documentclass[10pt,letterpaper]{article}
\usepackage{cornell-notes}

\title{Clause 33.05.07.05: Partial Specialization For Pointers}
\author{}
\date{}
\setDocTitle{Clause 33.05.07.05: Partial Specialization For Pointers}
\setDocSubtitle{C++ 2024}
\setDocOwner{C++ 2024 Cornell Notes}
\setDocDate{\today}
\setCornellCollection{C++ 2024 Cornell Notes}
\setCornellUnitType{Clause}
\setCornellUnitNumber{33.05.07.05}
\setCornellUnitTitle{Partial Specialization For Pointers}
\hypersetup{pdftitle={Clause 33.05.07.05: Partial Specialization For Pointers},pdfsubject={Cornell notes},pdfauthor={}}

\begin{document}
\maketitle
\begin{CornellOverviewBox}{Overview}
\textbf{Classification:} Standard library - Normative\\[2pt]
\textbf{Learning objective:} Understand the rules, terminology, and semantic consequences associated with Partial specialization for pointers.
\end{CornellOverviewBox}\begin{CornellSummaryBox}{Summary}
{\sffamily\bfseries\color{Accent} Summary}\par\vspace{3pt}Section 33.5.7.5 focuses on Partial specialization for pointers. Apply the specified interface together with its mandates, effects, result conditions, exception behavior, and complexity constraints. When applying it, preserve the distinction between mandatory requirements, recommendations, permissions, and behavior deliberately left undefined, unspecified, or implementation-defined.
\end{CornellSummaryBox}

\end{document}
